% from ega4-18-fr.tex
\oldpage[IV-4]{147}

  \item[(iii)] ${}^{hs}\!A_{(i)}$ est un $A$-module fidèlement plat, et si $\mathfrak m$ est l'idéal maximal de $A$, $\mathfrak m.{}^{hs}\!A_{(i)}$ est l'idéal maximal de ${}^{hs}\!A_{(i)}$, et $k({}^{hs}\!A_{(i)})$ est une clôture séparable de $k$.
  \item[(iv)] Pour que ${}^{hs}\!A_{(i)}$ soit noethérien, il faut et il suffit que $A$ le soit.
  \item[(v)] Si $A$ est strictement local (donc $\Omega=k$), on a ${}^{hs}\!A\xrightarrow{\sim}A$.
  \item[(vi)] Si $i':k\to\Omega'$ est un second homomorphisme de $k$ dans une clôture séparable de $k$, pour tout $k$-isomorphisme $\sigma:\Omega\xrightarrow{\sim}\Omega'$, le $A$-homomorphisme correspondant (par (ii)) $v_\sigma:{}^{hs}\!A_{(i)}\to{}^{hs}\!A_{(i')}$ est un isomorphisme.
\end{enumerate}
\end{proposition}

On a déjà vu ci-dessus que ${}^{hs}\!A_{(i)}$ est un anneau local et que $w$ est un homomorphisme local; le fait que ${}^{hs}\!A_{(i)}$ soit hensélien (et par suite strictement local) se démontre comme dans (18.6.6). Les assertions de (iii) découlent encore de ($0_{\mathrm{III}}$, 10.3.1.3), ainsi que la suffisance de la condition (iv); la nécessité de cette condition résulte de ce que ${}^{hs}\!A_{(i)}$ est un $A$-module fidèlement plat ($0_{\mathrm I}$, 6.5.2).

Pour prouver (ii), remarquons, avec les notations de (18.8.6), que pour tout $\lambda\in L$, il existe un et un seul $A$-homomorphisme local $v_\lambda:A_\lambda\to B$, tel que le composé $k(A_\lambda)\xrightarrow{\omega_\lambda}\Omega\xrightarrow{\psi}k(B)$ soit déduit de $v_\lambda$ par passage aux quotients, en vertu de (18.8.4) et de l'hypothèse que $B$ est strictement local. L'unicité de $v_\lambda$ entraîne en outre que les $v_\lambda$ forment un système inductif, d'où l'existence de l'homomorphisme $v$ répondant aux conditions de (ii); son unicité résulte de (17.4.9). L'unicité de $v$ entraîne aussitôt l'assertion (vi), en considérant les composés $v_\sigma\circ v_{\sigma^{-1}}$ et $v_{\sigma^{-1}}\circ v_\sigma$. Enfin, si $A$ est strictement local, et $A'=B_{\mathfrak n}$ une $A$-algèbre locale essentiellement étale, où $B$ est une $A$-algèbre étale, il résulte de (18.8.1) qu'il existe une $\Spec(A)$-section de $\Spec(B)$ prenant la valeur $\mathfrak n$ au point fermé de $\Spec(A)$, donc on déduit de (17.4.1) que l'homomorphisme $A\to A'$ est bijectif, ce qui prouve (v).

\begin{env}[18.8.8.1]
\label{IV.18.8.8.1}
Soit $\mathcal C$ la catégorie des anneaux strictement locaux, avec pour morphismes les homomorphismes locaux; pour tout $B\in\mathcal C$, désignons par $F(B)$ l'ensemble des couples $(u,\psi)$ formés d'un homomorphisme local $u:A\to B$ et d'un $k$-isomorphisme $\psi:\Omega\to k(B)$ tels que le composé $k(A)=k\xrightarrow{i}\Omega\xrightarrow{\psi}k(B)$ soit déduit de $u$ par passage aux quotients. On peut dire que l'objet ${}^{hs}\!A_{(i)}$ \emph{représente} le foncteur covariant $F:\mathcal C\to\mathbf{Ens}$ ($0_{\mathrm{III}}$, 8.1.11).

On remarquera qu'en vertu de (18.8.8, (vi)) les hensélisés stricts ${}^{hs}\!A_{(i)}$ sont tous $A$-isomorphes (pour les divers $k$-homomorphismes $i$ de $k$ dans des clôtures séparables de $k$), mais pour $i$ et $i'$ donnés, il y a en général une infinité de $A$-isomorphismes ${}^{hs}\!A_{(i)}\xrightarrow{\sim}{}^{hs}\!A_{(i')}$. De façon précise, le groupe des $A$-automorphismes de ${}^{hs}\!A_{(i)}$ est isomorphe au groupe de Galois de $\Omega$ sur $k$.
\end{env}

\begin{env}[18.8.9]
\EGAReaderTarget{ega4.c04.statement.18-8-9}{IV.18.8.9}
Soient maintenant $A$ un anneau \emph{semi-local}, $\mathfrak m_j$ ($1\leq j\leq r$) ses idéaux maximaux, et pour tout indice $j$, considérons un homomorphisme $i_j:k(A_{\mathfrak m_j})\to\Omega_j$ de $k(A_{\mathfrak m_j})$ dans une clôture séparable de ce corps. On appelle \emph{hensélisé strict} de $A$ relatif aux $i_j$ et l'on note ${}^{hs}\!A$ (lorsqu'il n'en résulte pas de confusion) le produit $\prod_{j=1}^r({}^{hs}(A_{\mathfrak m_j}))$ des hensélisés stricts des anneaux locaux $A_{\mathfrak m_j}$ relativement aux homomorphismes $i_j$.

